% ------------------------------------------------------------
% CHAPTER 24
% ------------------------------------------------------------

\chapter{Semantic Quantization and Symbolic Collapse}

The classical view of quantization presents it as a rounding operation that maps a continuous value to a discrete one. In engineering practice, this mapping is regarded as an unavoidable imperfection, a necessary concession to finite precision. Yet, when one examines the behavior of delta–sigma modulators, especially from the perspective of RSVP field theory and active inference, a deeper structure becomes apparent. Quantization is not merely rounding. It is a symbolic projection, a collapse of a continuous field onto a discrete manifold that carries semantic structure. This chapter develops the theory of semantic quantization, in which quantization is interpreted as a form of symbolic collapse analogous to representational discretization in cognitive systems.

The key insight is that the quantizer implements a mapping
\[
Q : \mathbb{R} \to \mathcal{A},
\]
where \(\mathcal{A}\) is a discrete set of admissible symbols. In standard quantization theory, \(\mathcal{A}\) is typically a uniformly spaced lattice. However, nothing in the mathematics of delta–sigma loops requires \(\mathcal{A}\) to be uniform, Euclidean, or even numerical. From the RSVP perspective, \(\mathcal{A}\) is a semantic manifold: a space of discrete representational states that correspond to different symbolic interpretations of the field. When the quantizer collapses a continuous state \(x\) to a symbol \(a \in \mathcal{A}\), it performs a selection among semantic categories, each representing an equivalence class of continuous configurations. The quantizer therefore partitions the continuous field into semantic regions, and these regions play the same role as receptive fields in perception.

To formalize this, consider the continuous input \(x\) as an element of a differentiable manifold \(M\). The quantizer implements a map
\[
Q : M \to \mathcal{A}.
\]
One may define the preimage of a symbol \(a\),
\[
U_{a} = Q^{-1}(a),
\]
which is an open or piecewise-open region of the manifold representing the semantic class corresponding to \(a\). These regions need not be convex or regular; in many systems, they are defined by topological constraints or representational equivalence classes. In cognitive systems, such regions correspond to perceptual categories, and the quantizer behaves analogously to a perceptual collapse.

This interpretation is strengthened by the observation that quantization error,
\[
e[n] = x[n] - Q(x[n]),
\]
is not merely numerical noise but semantic deviation. It measures the extent to which the continuous field lies away from the center of the semantic region in which it collapses. In cognitive systems, such deviations correspond to perceptual distortions or classification uncertainties. In delta–sigma modulation, these deviations act as entropy injections that must be compensated by the loop to restore representational fidelity.

The delta–sigma loop provides a powerful mechanism for handling semantic collapse. The feedback signal 
\[
u[n] - y[n]
\]
represents the discrepancy between the world and its symbolic representation. This discrepancy is not simply an error in the engineering sense but a representational tension between continuous reality and discrete symbolic form. The loop filter transports this tension through a vector field, shaping it into high-frequency modes where its semantic content becomes negligible. In this way, delta–sigma modulation implements a form of symbolic consistency enforcement: it maintains coherent representation by routing semantic deviations into irrelevant modes while preserving low-frequency structure.

The connection to derived geometry becomes evident once we interpret quantization as a derived fiber product. In derived algebraic geometry, the fiber product of two spaces along a map identifies points that satisfy a constraint. In quantization, the continuous space \(M\) and the discrete semantic space \(\mathcal{A}\) are related by the collapse map \(Q\). For each symbol \(a\), the derived fiber product
\[
M \times_{\mathcal{A}} \{a\}
\]
captures the infinitesimal structure of the collapse, including higher-order interactions and degeneracies. The quantizer thus performs a homotopy-corrected projection. While the naive collapse loses information, the derived collapse preserves certain higher-order structures that can be recovered through feedback and integration. This explains why delta–sigma modulators can reconstruct the continuous field so accurately despite the extreme loss of information at each quantization step.

To make this precise, consider the behavior of the error field \(e[n]\) under the loop dynamics. The recurrence relation for the state variable,
\[
x[n+1] = f(x[n]) + g(u[n] - y[n]),
\]
propagates the error through the system. Because the quantizer resets the symbolic state abruptly, the error behaves like a cochain whose differential is shaped by the loop filter. The loop therefore implements a discrete version of homotopy lifting: it reconstructs the path of the continuous field in the original manifold by following the high-order structure of the error, even though the observed symbols lie in the discrete semantic space.

This behavior mirrors perceptual inference in biological or cognitive systems. When sensory input is quantized into perceptual categories, the underlying continuous field must be inferred. The brain performs this inference through a combination of relaxation, prediction, and correction, which closely resembles the delta–sigma loop. The semantic category provides a coarse symbol, and internal dynamics reconstruct the fine structure. This reconstruction is effective because the continuous signal evolves on a smooth manifold with bounded curvature. Likewise, delta–sigma modulators can reconstruct arbitrary continuous signals precisely because they assume low curvature in the underlying field and correct deviations through feedback.

Dither provides an illuminating perspective on semantic quantization. Dither is traditionally understood as noise added to decorrelate quantization error. In the semantic view, dither injects controlled entropy into the collapse process, ensuring that the mapping from the continuous field to the semantic manifold remains unbiased. In cognitive systems, a similar mechanism exists in the form of exploratory noise, which prevents perceptual collapse from becoming overly deterministic and preserves sensitivity to subtle distinctions. From the RSVP perspective, dither smooths the entropy field \(S\), reducing curvature discontinuities and ensuring that the vector field can route entropy into high-frequency modes without destabilizing the system.

Quantization error, therefore, becomes a semantic measure. Its dynamics encode the mismatch between the manifold structure of the field and the discrete symbolic manifold \(\mathcal{A}\). The NTF of a delta–sigma modulator is not simply a spectral operator but a semantic smoothing operator, transporting representational deviations into modes where they do not affect meaning. In cognitive terms, delta–sigma modulators suppress irrelevant discrepancies and preserve coherent semantic structure.

Ultimately, semantic quantization reveals that delta–sigma modulation is a symbolic computation engine. It continually collapses continuous fields into discrete categories and reconstructs the underlying structure through geometric flows that preserve coherence. The next chapter deepens this interpretation by analyzing the geometry of error on curved manifolds, providing a unified geometric framework for delta–sigma stability, accuracy, and semantic fidelity.


